\section*{Chapter \ref{ch:g10:drainage-basins} --- Drainage Basins, Rivers and Floods}

\begin{solution}{exo:g10:drainage-basins:1}
The \omterm{def:g10:drainage-basins:drainage-basin}{drainage basin} of a point of a river is the whole area of land whose
runoff flows to the river upstream of that point and passes it. The \omterm{def:g10:drainage-basins:drainage-basin}{drainage
divide} is the line, along the crests, that separates two neighbouring basins.
Two drops falling a metre apart on either side of a divide run down opposite
slopes, into two different basins: their rivers can reach different seas.
\end{solution}

\begin{solution}{exo:g10:drainage-basins:2}
The \omterm{def:g10:drainage-basins:discharge}{discharge} is the volume of water passing through a cross-section of the
river each second, in \unit{m^3/s}. \qty{1}{m^3/s} is \qty{1000}{L} each
second; in a day, $1000 \times 86\,400 = 86.4$ million litres.
\end{solution}

\begin{solution}{exo:g10:drainage-basins:3}
Before the flood the river carried about \qty{330}{m^3/s} (between about
\num{240} and \qty{390}{m^3/s}). The peak was \qty{9231}{m^3/s} on
21 January. The rise began on 18 January (\qty{362}{m^3/s}): it took three
days.
\end{solution}

\begin{solution}{exo:g10:drainage-basins:4}
$1/50 = 0.02$, \qty{2}{\percent}. No: the years are independent; the chance
is \qty{2}{\percent} again this year, whatever happened last year.
\end{solution}

\begin{solution}{exo:g10:drainage-basins:5}
One cubic foot is $0.3048^3 \approx \qty{0.02832}{m^3}$, so
$593\,000 \times 0.02832 \approx \qty{16800}{m^3/s}$. The Amazon carries
$206\,000/16\,800 \approx 12$ times more.
\end{solution}

\begin{solution}{exo:g10:drainage-basins:6}
Strip areas: $1.5 \times 0.3 = 0.45$, $0.90$, $1.05$ and \qty{0.60}{m^2}:
$A = \qty{3.0}{m^2}$. \omterm{def:g10:drainage-basins:discharge}{Discharges}: $0.2 \times 0.45 = 0.09$, $0.5 \times 0.90
= 0.45$, $0.6 \times 1.05 = 0.63$, $0.3 \times 0.60 = 0.18$: $Q =
\qty{1.35}{m^3/s}$. Mean velocity $1.35/3.0 = \qty{0.45}{m/s}$.
\end{solution}

\begin{solution}{exo:g10:drainage-basins:7}
$66 \times 10^6 \times 0.020 = \qty{1.32e6}{m^3}$. Two days are
$2 \times 86\,400 = \qty{172800}{s}$: $1.32 \times 10^6/172\,800 \approx
\qty{7.6}{m^3/s}$ on average.
\end{solution}

\begin{solution}{exo:g10:drainage-basins:8}
Ranked: 120, 88, 65, 59, 51, 47, 42, 38, 35. With $n = 9$, $T = 10/m$: the
largest (120) has $T = 10$ years, the second (88) 5 years, the third (65)
about 3.3 years. The median, 51, has rank 5 and $T = 2$ years: it is equalled
or exceeded one year in two.
\end{solution}

\begin{solution}{exo:g10:drainage-basins:9}
The town's \omterm{def:g10:drainage-basins:hydrograph}{hydrograph} rises sooner and much higher, and falls quickly: on
steep slopes, roofs and roads, little water soaks in and the runoff reaches
the stream fast and all at once; its \omterm{def:g10:drainage-basins:hydrograph}{lag time} is short. The forested basin's
\omterm{def:g10:drainage-basins:hydrograph}{hydrograph} rises later, to a lower, broader peak: the soil and the forest
take in much of the rain and release it slowly, and the water flows slowly on
gentle slopes. Its \omterm{def:g10:drainage-basins:hydrograph}{baseflow} stays higher after the storm.
\end{solution}

\begin{solution}{exo:g10:drainage-basins:10}
$9231 \times 86\,400 \approx \qty{7.98e8}{m^3}$, spread over
$29\,900 \times 10^6 = \qty{2.99e10}{m^2}$: about \qty{0.027}{m},
\qty{27}{mm}. The \omterm{def:g10:drainage-basins:hydrograph}{baseflow} gives $330 \times 86\,400/2.99 \times 10^{10}
\approx \qty{0.001}{m}$, about \qty{1}{mm} a day: the peak day carried about
28 times more.
\end{solution}

\begin{solution}{exo:g10:drainage-basins:11}
$P = 1 - 0.9^{10} \approx 1 - 0.349 = 0.651$: about \qty{65}{\percent}. It
does not come at all with probability $0.9^{10} \approx 0.35$, about
\qty{35}{\percent}.
\end{solution}

\begin{solution}{exo:g10:drainage-basins:12}
The \omterm{def:g10:drainage-basins:hydrograph}{hydrograph} is a triangle three days long and \qty{1200}{m^3/s} high: its
area is $\tfrac12 \times 1200 \times 3 \times 86\,400 \approx
\qty{1.56e8}{m^3}$, \num{156} million cubic metres. The part above
\qty{800}{m^3/s} is a smaller triangle of the same shape, three times
smaller in height ($400$ instead of $1200$) and in duration (one day instead
of three): $\tfrac12 \times 400 \times 86\,400 \approx \qty{1.73e7}{m^3}$,
\num{17} million cubic metres. That is less than the \num{60} million the
reservoir can store: yes.
\end{solution}

\begin{solution}{exo:g10:drainage-basins:13}
Full to its banks, the river is $800/(100 \times 2) = \qty{4}{m}$ deep.
(1) The floodplain carries \qty{400}{m^3/s} over \qty{900}{m} at
\qty{0.5}{m/s}: its depth is $400/(900 \times 0.5) \approx \qty{0.9}{m}$.
(2) Between the levees: $1200/(100 \times 2) = \qty{6}{m}$ deep, \qty{2}{m}
above its banks, more than twice the depth on the floodplain without
levees.
\end{solution}

\begin{solution}{exo:g10:drainage-basins:14}
We want $(1 - 1/T)^{30} = 0.5$, so $1 - 1/T = 0.5^{1/30} \approx 0.97716$,
$1/T \approx 0.02284$ and $T \approx 44$ years.
\end{solution}

\begin{solution}{exo:g10:drainage-basins:15}
Each square sends its water to its lowest neighbour (rows from north to south,
columns from west to east):
\begin{center}
\begin{tabular}{|c|c|c|c|}
\hline
outlet & west & west & west \\ \hline
north & west & west & south \\ \hline
north & north & south & south \\ \hline
north & east & east & outlet \\ \hline
\end{tabular}
\end{center}
Following the arrows, the squares 47, 39, 44, 37, 46 and 30 drain to the
south-east outlet: 6 squares, \qty{0.06}{km^2}. The other 10 squares drain
to the north-west outlet: \qty{0.10}{km^2}. The divide runs from the east
edge between 50 and 47, then between 48 and 47, between 48 and 44, between
49 and 44, between 49 and 46, and between 51 and 46 to the south edge.
\end{solution}

\begin{solution}{pb:g10:drainage-basins:1}
\textbf{1.} $11\,560 \times 1.609^2 \approx \qty{29900}{km^2}$.

\textbf{2.} The \omterm{def:g10:drainage-basins:drainage-basin}{drainage divide} of the basin upstream of the station.

\textbf{3.} $484\,000 \times 0.3048^3 \approx \qty{13700}{m^3/s}$.

\textbf{4.} $13\,705/3341 \approx 4.1$ times.

\textbf{5.} $A = 400 \times 8 = \qty{3200}{m^2}$ and $v = Q/A = 13\,705/3200
\approx \qty{4.3}{m/s}$.

\textbf{6.} About \qty{330}{m^3/s}.

\textbf{7.} \qty{9231}{m^3/s}, on 21 January.

\textbf{8.} Three days.

\textbf{9.} $9231 \times 86\,400/(29\,900 \times 10^6) \approx
\qty{0.027}{m}$: about \qty{27}{mm}.

\textbf{10.} New water reached the river: another fall of rain on the basin
(or melting snow) sent a second, smaller pulse of runoff.

\textbf{11.} Seventeen days after the peak. The rise comes from the runoff of
the slopes near the streams, which arrives quickly and all at once; on the
falling limb the river still receives the water of the farthest parts of the
basin and the water that drains slowly out of the soil.

\textbf{12.} $T = n + 1 = 96$ years.

\textbf{13.} $96/2 = 48$ years.

\textbf{14.} $96/3 = 32$ years.

\textbf{15.} $96/4 = 24$ years and $96/5 = 19.2$ years. Not much: either way
this \omterm{def:g10:drainage-basins:discharge}{discharge} comes back about once in twenty years.

\textbf{16.} Five years (1936, 1937, 1943, 1972 and 1996): $5/95 \approx
0.053$, about \qty{5}{\percent} a year, a \omterm{def:g10:drainage-basins:return-period}{return period} of about
$96/5 \approx 19$ years.

\textbf{17.} $1/100 = 0.01$, \qty{1}{\percent}.

\textbf{18.} $1 - 0.01 = 0.99$.

\textbf{19.} $0.99^{30} \approx 0.740$.

\textbf{20.} $1 - 0.740 = 0.26$: \textbf{the chance of at least one
100-year flood in 30 years is about \qty{26}{\percent}}, one in four.
\end{solution}
